\subsection{ms-evaluaciones --- Tareas y entregas}
\label{sec:ms-evaluaciones}
\textit{Responsable: Edinson (EVA-01). Repositorio: \url{https://github.com/nubeduca/nubeduca-ms-evaluaciones}.}

\pendiente{Dos o tres párrafos: qué resuelve este microservicio, decisiones de diseño y qué roles (STUDENT, TEACHER, ADMIN) usan cada operación.}

\subsubsection{Endpoints}
\begin{table}[htbp]
\centering\footnotesize
\begin{tabularx}{\textwidth}{lXll}
\toprule
Método & Ruta & Roles & Emite \\
\midrule
GET    & \texttt{/health} & pública & --- \\
GET    & \texttt{/courses/\{courseId\}/assignments} & todos & --- \\
POST   & \texttt{/courses/\{courseId\}/assignments} & TEACHER & --- \\
GET    & \texttt{/assignments/\{assignmentId\}} & todos & --- \\
PATCH  & \texttt{/assignments/\{assignmentId\}} & TEACHER & --- \\
POST   & \texttt{/assignments/\{assignmentId\}/publish} & TEACHER & \texttt{assignment.published} \\
GET    & \texttt{/assignments/\{assignmentId\}/submissions} & STUDENT, TEACHER & --- \\
POST   & \texttt{/assignments/\{assignmentId\}/submissions} & STUDENT & \texttt{submission.submitted} \\
GET    & \texttt{/submissions/\{submissionId\}} & STUDENT, TEACHER & --- \\
POST   & \texttt{/submissions/\{submissionId\}/grade} & TEACHER & \texttt{submission.graded} \\
\bottomrule
\end{tabularx}
\caption{Endpoints de ms-evaluaciones (\texttt{apis/ms-evaluaciones.openapi.yaml} de \texttt{nubeduca-contratos}); \textit{todos} = STUDENT, TEACHER y ADMIN.}
\end{table}

\subsubsection{Eventos}
\textbf{Publica:}
\begin{itemize}
  \item \texttt{lms.assessment.assignment.published.v1}
  \item \texttt{lms.assessment.submission.submitted.v1}
  \item \texttt{lms.assessment.submission.graded.v1}
\end{itemize}

\textbf{Consume:}
\begin{itemize}
  \item \texttt{lms.courses.enrollment.created.v1}
  \item \texttt{lms.courses.enrollment.dropped.v1}
  \item \texttt{lms.content.file.scanned.v1}
\end{itemize}


\subsubsection{Estructura JSON de la tabla \texttt{<stage>-t\_evaluaciones}}
Clave de partición \texttt{tenant\_id}, clave de ordenamiento \texttt{sk}.

\begin{table}[htbp]
\centering\small
\begin{tabularx}{\textwidth}{>{\raggedright\arraybackslash}p{0.42\textwidth}X}
\toprule
\texttt{sk} & Qué es \\
\midrule
\texttt{ASSIGN\#<assignmentId>} & Tarea \\
\texttt{COURSEASSIGN\#<courseId>\#<dueAt>\#<assignmentId>} & Índice: tareas de un curso por fecha de entrega \\
\texttt{SUBMIT\#<submissionId>} & Entrega \\
\texttt{ASSIGNSUBMIT\#<assignmentId>\#<studentId>\#<submissionId>} & Índice: entregas de una tarea (y de un alumno en esa tarea) \\
\texttt{ATTEMPTS\#<assignmentId>\#<studentId>} & Contador de entregas del alumno en la tarea (\texttt{total}) \\
\texttt{PROY\#ENROLL\#<courseId>\#<userId>} & Matrícula, proyección de \texttt{enrollment.created} \\
\texttt{PROY\#FILE\#<fileId>} & Archivo escaneado, proyección de \texttt{file.scanned} \\
\texttt{EVENTO\#}, \texttt{INBOX\#}, \texttt{AGG\#} & Outbox, inbox y versión de agregados de la plantilla \\
\bottomrule
\end{tabularx}
\caption{Prefijos de \texttt{sk} en la tabla de ms-evaluaciones.}
\end{table}

\begin{lstlisting}
{
  "tenant_id": "utec",
  "sk": "ASSIGN#asg_01JCF4SQ7XAMY8HTMRFHAWZ73F",
  "entidad": "ASSIGNMENT",
  "assignmentId": "asg_01JCF4SQ7XAMY8HTMRFHAWZ73F",
  "courseId": "crs_01JCQRZHQT9F4WY2JYPDV546K7",
  "title": "Laboratorio 4: SQS y DLQ",
  "instructions": "Sube tu informe en PDF.",
  "dueAt": "2026-11-20T23:59:00Z",
  "maxScore": 20,
  "allowLate": false,
  "status": "PUBLISHED",
  "ownerId": "c62d9930-b473-4b1d-adec-72fdda7cf599",
  "createdAt": "2026-11-11T15:00:00Z",
  "version": 2
}
\end{lstlisting}

\begin{lstlisting}
{
  "tenant_id": "utec",
  "sk": "SUBMIT#sbm_01JCY7SYF5X0ZHRJ7KQSVX4BP0",
  "entidad": "SUBMISSION",
  "submissionId": "sbm_01JCY7SYF5X0ZHRJ7KQSVX4BP0",
  "assignmentId": "asg_01JCF4SQ7XAMY8HTMRFHAWZ73F",
  "courseId": "crs_01JCQRZHQT9F4WY2JYPDV546K7",
  "studentId": "a18577e2-e005-4a49-a29e-755dbf08f5b6",
  "attempt": 1,
  "fileIds": ["fil_01JC8P2K48BKYRT5ATH3DZDA66"],
  "comment": "Adjunto mi informe.",
  "late": false,
  "status": "GRADED",
  "submittedAt": "2026-11-12T15:22:00Z",
  "score": 17,
  "feedback": "Buen trabajo.",
  "graderId": "c62d9930-b473-4b1d-adec-72fdda7cf599",
  "gradedAt": "2026-11-13T15:23:00Z",
  "version": 2
}
\end{lstlisting}

\begin{lstlisting}
{
  "tenant_id": "utec",
  "sk": "PROY#ENROLL#crs_01JCQRZHQT9F4WY2JYPDV546K7#a18577e2-e005-4a49-a29e-755dbf08f5b6",
  "entidad": "PROYECCION",
  "enrollmentId": "enr_01JC5F0KPBVF7HK1VV8H23JGMW",
  "courseId": "crs_01JCQRZHQT9F4WY2JYPDV546K7",
  "userId": "a18577e2-e005-4a49-a29e-755dbf08f5b6",
  "role": "STUDENT",
  "recibidoDe": "lms.courses",
  "version": 1
}
\end{lstlisting}

\begin{lstlisting}
{
  "tenant_id": "utec",
  "sk": "PROY#FILE#fil_01JC8P2K48BKYRT5ATH3DZDA66",
  "entidad": "PROYECCION",
  "fileId": "fil_01JC8P2K48BKYRT5ATH3DZDA66",
  "result": "CLEAN",
  "purpose": "SUBMISSION",
  "ownerId": "a18577e2-e005-4a49-a29e-755dbf08f5b6",
  "recibidoDe": "lms.content",
  "version": 1
}
\end{lstlisting}

Formato de los ítems documentado en el README de \texttt{nubeduca-ms-evaluaciones}.

\subsubsection{Datos de prueba}
\pendiente{Evidencia de los 30\,000 registros (UTEC 15\,000, UPC 10\,000, PUCP 5\,000): salida de \texttt{poblar\_datos.py} o un conteo por tenant con \texttt{aws dynamodb query ... --select COUNT}.}

\subsubsection{Resultados y evidencias}
\evidencia{ms-evaluaciones-evidencia.png}{ms-evaluaciones: prueba funcional (por ejemplo, una llamada en Postman o una captura del frontend).}{fig:ms-evaluaciones}
